% !TeX root = ../papers/nc-p3.tex
\section{階数三の対称核の二項四次部分}
\label{sec:cone-kernel-binary-transfer}
\begin{theorem}\label{thm:cone-kernel-binary-transfer}
四辺が有限な標準有限Frobenius窓が片側対称で、対称二乗に階数3の核元があるとする。
対称な二位置の点座標を $E=k e\oplus S^2 B$, $\dim B=2$ とし、
核元を二項二次形式の判別式に対応させる。
前節と同じ六次元への隣接分解 $\beta,\gamma$ に対し、一変数の二つの鉛筆
\[
 P(a)(b\otimes f)=\beta(ab,f),\qquad Q(a)(b\otimes g)=\gamma(ab,g)
 \quad(a,b\in B)
\]
を作る。双方が一般に全射なら窓は両対称である。
少なくとも一方が一般に非全射なら元代数は指数成長する。
特に長さ1の点を持つ反対階数9の全有限点窓で全次数の正則性は有限判定できる。
\end{theorem}
核元を消すことは、二つの $S^2 B$ への制限で四つの二項変数の全てが対称であることに等しい。
実際 $S^2(S^2 B)\to S^4 B$ の積の核は判別式の一次元である。
また関係に純テンソル $u f$ があれば、隣り合う転置窓で $u f u$ が三次重なりに入り、
その三テンソルに対応する非零点で辺の評価階数が1となり有限性に反する。
従って $S^2 B$ 上の $\beta,\gamma$ の値は各々六次元標的を張る。
実際、張らなければそれを消す汎関数の関係は $e$ 成分だけにあり純関係となる。

\begin{proof}[両方が全射となる方向]
一般の二項基底 $x,y$ で $P(x),Q(x)$ は全射とする。四二項位置の対称性から
前節の二行列分解と同じく $P(y)=TP(x),\ Q(y)=T^tQ(x)$ である。
従って $x^2,xy,y^2$ で評価した写像は、それぞれ
$U,TU,T^2U$ および $V,T^tV,(T^t)^2 V$ である。
ここで $U:F\to H,\ V:G\to H^*$ は単射である。実際、$[x^2]$ は有限点の外に選べる。
さらに $U,TU$ の像は合わせて $H$ を張り、$V,T^tV$ の像は $H^*$ を張る。
$U',V'$ を $e$ での評価とすれば、偶位置対称性により
\[
 (T^jU')^t V=(T^j U)^t V'\qquad(j=0,1,2).
\]
これらの式を使って $S(Uf+TUf')=U'f+TU'f'$ と置く。
左辺の引数が零なら右辺は $V,T^t V$ を消すので零であり、$S$ は良定義である。
また $S^t V=V',\ S^t T^t V=T^t V'$ も同じ式による。
$C=[S,T]$ とすると $CU=0,\ C^t V=0$ である。
各 $x^2,xy,y^2,e$ での評価は、以後 $M U,M^t V$ と書ける。
ここで $M$ は順に $I,T,T^2,S$ である。

$U$ の像の余次元から $\operatorname{rank} C\le2$ だが、実は等しい。
$h\in\ker C^t$ ごとに $M\mapsto M^t h$ は $P(\beta)$ の対称三次重なりを与える。
実際 $h^t[M,N]U=0$ を調べればよく、唯一残る $[S,T^2]=CT+TC$ でも $h^t C=0$, $CU=0$ である
（ここで $h^t$ は行汎関数を表す）。$I$ での評価によりこれらは独立で、
三次重なりは四次元である。従って
\[
 \operatorname{rank} C=2,\qquad U(F)=\ker C,\quad V(G)=\ker C^t .
\]
各固有値 $\lambda$ で $T$ の固有空間と $U(F)$ の交叉は零である。
交叉元 $Uf\ne0$ があれば、$p(\lambda)=0$ を満たす全ての二次以下の多項式で
$p(T)Uf=0$ となり、偶文字の射影直線が辺に射影される。転置についても同様である。
$C=A D$ を $A:k^2\hookrightarrow H,\ D:H\twoheadrightarrow k^2$ と階数分解する。
同じ固有値の左右固有ベクトル $s,r$ には $sCr=s[S,T]r=0$ が成り立つ。
もし固有空間が二次元以上なら、$D$ のその上での単射性と左固有空間上での $A^t$ の単射性により
$sA,Dr$ はそれぞれ二次元全部を走り、矛盾する。従って $T$ は巡回である。

六次の最小多項式 $p$ により $H=k[t]/p$ とし、最高次係数汎関数 $\ell$ の完全積対合を使う。
$C(s)=u_0\ell(v_1 s)-u_1\ell(v_0s)$ と書く。
$\operatorname{tr}(T^j [S,T])=0$ より、全 $j$ に対して $\ell(t^j(u_0v_1-u_1v_0))=0$ であり、
$u_0v_1=u_1v_0$ となる。各根で $u$ の二元がともに零なら、その根での評価が左固有汎関数で
$C$ の像を消してしまう。$v$ についても、二元がともに零ならその局所因子のsocleは
$T$ の右固有空間で $C$ に殺される。いずれも先の交叉零に反する。
従って $v_i=h u_i$ ($i=0,1$) となる単元 $h\in H^\times$ がある。
完全対称対合 $\ell(h a b)$ では $T$ は自己随伴、$C$ は反自己随伴である。
この随伴を $\dagger$ とすると $[S^\dagger,T]=C$、ゆえに $S-S^\dagger$ は巡回な $T$ と交換し、
$T$ の多項式で自己随伴となる。反自己随伴でもあるため $S=S^\dagger$ である。
最後に $\ker C$ と $\ker C^t$ はこの対合で対応し、$U,V$ で二つの奇位置を同定できる。
全 $M$ が自己随伴なので得られる窓はGram表示の両対称窓である。
\end{proof}

\begin{proof}[非全射の鉛筆]
必要なら奇位置を交換し、$P$ が一般に非全射とする。一般階数は5である。
実際4以下なら各 $a$ に対して $\mathbb P B\times\mathbb P F$ 上の高々四式の共通零点があり、
有限個の辺点では全 $a$ を覆えない（二項二次形式の一次因子の個数は有限）。
原始左核 $l(a)$ の次数は五次以下である。$\beta(ab,-)$ の一般階数は4で、
$ab$ の評価の像は $l(a),l(b)$ を消す。この二元が一般に比例すれば $l$ は定数となり
六次元を張ることに反する。最大四次小行列式の相補座標の次数は $(4,4)$ で、
退化が有限なので共通因子を持たない。従ってこれらは $l(a)\wedge l(b)$ の原始化である。
対角 $a=b$ が共通因子なので $\deg l\ge5$、ゆえに等しい。
共通右核を割った列数は5である。もし6以上なら、独立な入力から一般階数5の六列を選び、
随伴行列を五次の $l$ でくくると定数右核が生じて矛盾する。
残る六行五列の鉛筆の最大階数は全点で5である（最大五次小行列式は $l$ の非零定数倍）。
完全列 $5\mathcal O(-1)\to6\mathcal O\to\mathcal O(5)$ により
$l$ は有理正規五次曲線であり、$H=S^5B$ と取ると $l(a)$ の核は五次式のうち $a$ の倍数全体である。
従って $\beta(ab,f)$ は $a,b$ 両方の倍数である。普遍一次式の互いに素な二因子で割って
$\beta(z,f)=zf$ ($z\in S^2 B,\ f\in F=S^3 B$) を得る。
商の $F\to S^3 B$ は単射である。核があれば固定した奇点の相手が射影平面全部を含む。

$D:S^3 B\to S^5 B$ を $\beta(e,-)$ と書き、基底 $x,y$ で
\[
 K=\langle xD(yu)-yD(xu):u\in S^2 B\rangle\subset S^6 B
\]
と置く。$G$ は $S^7 B$ 上で $BK$ を消す全汎関数と同定できる。
実際、四二項位置の対称性と $zf$ の五変数の全対称性から、各 $g$ で
$w(z,f,z',g)=L_g(z f z')$ となる七次の汎関数がある。
これは五変数と四変数の置換が重なった二位置を介して七変数の対称群を生成することによる。
$e$ との対称性は $\gamma(e,g)(zf)=L_g(z D f)$ である。
$S^2 B\otimes S^3 B\to S^5 B$ の乗法核は $xv\otimes yu-yv\otimes xu$ ($v\in B,u\in S^2 B$)
で張られる。各全次数の単項式の隣り合う分解の差を取ればよい。
従って $\gamma(e,g)$ の存在・一意性は $L_g(BK)=0$ と同値である。
逆にその任意の $L$ は同じ式から全偶位置対称の三次重なりを定める（$e,e$ の条件は自動）。
三次重なりの次元は4であり、その中に元の $G$ が単射で入るから主張が従う。
特に $\dim BK=4$、$\dim K\ge2$ である。

一次縮約 $J:B\otimes G\to (S^6 B)^*$ を $J(a,g)(s)=L_g(as)$ とする。
全像は $K$ を消すので階数は $7-\dim K\le5$ である。
両端対称な $B\otimes G\otimes B$ の部分（次元12）において、
左から $J$ で縮約して零となる部分の次元は少なくとも $12-2\operatorname{rank}J\ge2$ である。
対称性から右の縮約でも零となる。
その非零元 $h$ を固定する。
偶位置を $S^2 B$ に制限した経路の双対テンソルを作る。
各内部の $E,G,E$ 三位置には $h$ を一つ置き、各 $E,F$ 辺には
二つの二項変数の角括弧を一つずつ掛ける。
例えば両端を $F$ とする長さ $4s+1$ の列を使えば、$s$ 個の三位置は互いに素で、
偶位置は一次と括弧の一次で丁度二次になる。
各 $F$ 位置は三次まで任意の残りの因子を掛けられる。
全 $F$ 辺で括弧により乗法は零で、$G$ 辺でも $J$ の零条件が、偶位置に
さらに掛けた一次因子によらず $\gamma(z,g)(s')=L_g(zs')$ を消す。
よって得たテンソルは全関係を消す。固定した括弧積と互いに異なる位置の $h$ の積は非零である。
各内部 $F$ の残る一次を独立に選んだテンソルは独立である
（$G$ の係数を評価して各 $h$ を非零の二項多項式にしてから整域で固定積を消去してよい）。
従って次元は少なくとも $2^{s-1}$ であり、指数成長する。
\end{proof}

\begin{corollary}\label{cor:rank-nine-simple-ternary-equivalence}
反対階数9の有限点グラフに長さ1の局所点があるとする。標準有限Frobenius検査を通った窓が
全次数で正則であることは、いずれかの点で中央の支持が3であることと同値である。
特に全正則例は三文字部分代数の既出の表示を持つ。
\end{corollary}
\begin{proof}
中支持3の十分性は定理\ref{thm:finite-cubic-separators}である。全中支持が4なら
定理\ref{thm:rank-nine-finite-symmetry}で片側対称な表示があり、核の階数は3か4である。
この二つの階数では前定理と定理\ref{thm:smooth-symmetric-kernel-all}から
両対称または指数成長である。両対称な標準AS窓の反対階数は6か8なので
（定理\ref{thm:bisymmetric-six-or-eight}）、正則例はない。
\end{proof}
